\documentclass[a4paper,12pt]{article}
\usepackage{amsmath, amssymb}
\usepackage{xcolor}
\usepackage[most]{tcolorbox}
\usepackage{geometry}

\geometry{left=2cm, right=2cm, top=2cm, bottom=2cm}

% Couleurs pour l'arrière-plan des box
\definecolor{SolutionColor}{HTML}{F0F8FF}
\definecolor{ExerciseColor}{HTML}{E6E6FA}

% Environnements stylisés
\newtcolorbox{solutionbox}[1]{
    colback=SolutionColor,
    colframe=blue!75!black,
    title=#1,
    fonttitle=\bfseries,
    breakable=true
}
\newtcolorbox{exercice}[1][]{
    colback=ExerciseColor,
    colframe=purple!75!black,
    title=#1,
    fonttitle=\bfseries
}

\begin{document}

%%%%%%%%%%%%%%%%%%%%%%%%%%%%%%%%%%%%%%%%%%%%%%%%%%%%%%%%%%%%%%
% Exercice 10
%%%%%%%%%%%%%%%%%%%%%%%%%%%%%%%%%%%%%%%%%%%%%%%%%%%%%%%%%%%%%%
\begin{exercice}[title=Exercice 10 $\star\star$ {[Calculer]}]
\begin{enumerate}
  \item Montrer que pour tous $x,\,y \in \mathbb{R}$ :
    $$
    x^3 - y^3 \;=\;(x - y)\,\bigl(x^2 + x\,y + y^2\bigr).
    $$
  \item Déterminer l'ensemble des nombres premiers de la forme $n^3 - 8$, où $n$ est un nombre entier.
  \item Montrer que l'entier $n^3+1$ n'est pas premier.
\end{enumerate}
\end{exercice}

\begin{solutionbox}{Solution de l'exercice 10}

\textbf{1) Factorisation de $x^3-y^3$:}

Pour $x,y\in\mathbb{R}$, on sait la factorisation classique :
\[
x^3 - y^3 = (x-y)\,(x^2 + x\,y + y^2).
\]
Cette identité se vérifie en développant le membre de droite.  
\[
(x-y)(x^2 + x\,y + y^2)=x^3 + x^2y + xy^2 - (x^2y +xy^2 +y^3)=x^3-y^3.
\]

\textbf{2) Nombres premiers de la forme $n^3-8$:}

Observons $n^3-8= (n-2)(n^2+2n+4)$ par la formule précédente (avec $x=n,y=2$).  
\[
n^3-8 = (n-2)\bigl(n^2 +2n +4\bigr).
\]

Pour que $n^3-8$ soit premier, il doit être >1 et ne posséder que 1 et lui-même comme diviseurs. Or on voit deux facteurs $(n-2)$ et $(n^2+2n+4)$.  
Les cas possibles sont :

- \textbf{Cas 1 :} $n-2=1$, ce qui donne $n=3$. Alors $n^2+2n+4= 9+6+4=19$.  
  Ainsi $n^3-8= (3^3)-8=27-8=19$, qui est bien premier.

- \textbf{Cas 2 :} $n^2+2n+4=1$.  
  Résolvons $n^2 +2n +4=1 \implies n^2+2n+3=0\implies (n+1)^2= -2,$ pas de solution réelle.  
  Aucun $n$ entier ne satisfait.

- \textbf{Cas 3 :} $n-2=-1$ ou $n^2+2n+4=-1$ ne permettent pas $n^3-8>1$, etc.

Donc la seule valeur $n$ pour laquelle $n^3-8$ est premier est $n=3$. Le nombre obtenu est $19$.

\textbf{Conclusion :} L’ensemble cherché est $\{\,19\}$ (issu de $n=3$).


\textbf{3) Montrer que $n^3+1$ n’est pas premier}

On veut prouver que pour tout entier $n\ge1,$ l’entier $n^3+1$ n’est pas un nombre premier.

\begin{itemize}
  \item \textit{Idée :} On peut factoriser $n^3+1$ comme $(n+1)(n^2-n+1)$, si on se rappelle la formule:
    \[
     x^3+1=(x+1)\bigl(x^2-x+1\bigr).
    \]
    Ici, $x=n$.
  
  \item Ainsi :
    \[
     n^3+1 = (n+1)(n^2-n+1).
    \]
    Pour $n\ge1$, on a $n+1\ge2$ et $n^2-n+1\ge 1$ ?? Plus précisément, si $n>1,$ on a $n+1\ge2$ et $n^2-n+1\ge 1$; en fait, pour $n=1,$ $n^2-n+1=1,$ mais $n+1=2$. Ça donne $1^3+1=2$, prime ou pas ? $2$ est bien un nombre premier, contradiction ?? On doit vérifier la condition. Peut-être la question suppose $n>1$.  

  \item Vérifions $n=1 : 1^3+1=2$, c’est un nombre premier. Éventuellement on comprend la question comme «pour $n\ge2$, prouver $n^3+1$ composite».  

    - \underline{Si $n=2$:} $n^3+1=8+1=9=3^2$, pas premier.
    - \underline{Si $n\ge2$:} $n+1\ge3$, $n^2-n+1\ge 3$ (on peut vérifier $n^2-n+1=(n-\tfrac12)^2 +\tfrac34>1$), donc on a un produit d'entiers >1, c’est donc composé.

\end{itemize}

\textbf{Conclusion :} Pour $n\ge2$, $n^3+1$ est décomposable en $(n+1)(n^2-n+1)$, chacun $\ge2$, donc $n^3+1$ n’est pas premier.

\end{solutionbox}

%%%%%%%%%%%%%%%%%%%%%%%%%%%%%%%%%%%%%%%%%%%%%%%%%%%%%%%%%%%%%%%%%%%%%%
% Exercice 12
%%%%%%%%%%%%%%%%%%%%%%%%%%%%%%%%%%%%%%%%%%%%%%%%%%%%%%%%%%%%%%%%%%%%%%
\begin{exercice}[title=Exercice 12 $\star\star$ -- [Chercher]]
Déterminer les trois plus petits entiers naturels $n$ tels que $n^2 -1$ soit le produit de trois nombres premiers distincts.
\end{exercice}

\begin{solutionbox}{Solution de l'exercice 12}
\textbf{Idée principale :} On veut $n^2 -1=(n-1)(n+1)$ soit le produit de trois nombres premiers distincts.  
Cela signifie $(n-1)$ et $(n+1)$ se factorisent en un total de trois facteurs premiers.  
On cherche les plus petits $n$.

\textbf{Approche :}  
\begin{enumerate}
  \item Noter que $n+1$ et $n-1$ sont deux entiers consécutifs pairs, si $n$ est pair. S’il est pair, alors $n-1$ et $n+1$ sont impairs. Etc.  
  \item Tester des petits $n$, factoriser $(n-1)(n+1)$ et voir si on obtient trois premiers distincts.

\item \textbf{Exemples} :
\[
\bullet\ n=2:\quad n^2-1=3,\ \text{qui n’a qu’un seul facteur premier.}
\]
\[
\bullet\ n=3:\quad n^2-1=8=2^3,\ \text{trois facteurs identiques (2,2,2) pas distincts.}
\]
\[
\bullet\ n=4:\quad n^2-1=15=3\times5,\ \text{deux facteurs distincts.}
\]
\[
\bullet\ n=5:\quad n^2-1=24=2^3\times3,\ \text{deux distincts, 2 et 3, mais un répété.}
\]
\[
\bullet\ n=6:\quad n^2-1=35=5\times7,\ \text{deux distincts.}
\]
\[
\bullet\ n=7:\quad n^2-1=48=2^4\times3,\ \text{deux distincts.}
\]
\[
\bullet\ n=8:\quad n^2-1=63=3^2\times7,\ \text{deux distincts.}
\]
\[
\bullet\ n=9:\quad n^2-1=80=2^4\times5,\ \text{deux distincts.}
\]
\[
\bullet\ n=10:\quad n^2-1=99=3^2\times11,\ \text{deux distincts.}
\]
\[
\bullet\ n=11:\quad n^2-1=120=2^3\times3\times5,\ \text{trois distincts 2,3,5.}
\]
Ça y est, $n^2-1=120=2^3\times3\times5$. On a trois distincts : (2,3,5).

Donc $n=11$ donne un premier exemple. On continue un peu :

\[
\bullet\ n=12:\quad n^2-1=143=11\times13,\ \text{deux distincts.}
\]
\[
\bullet\ n=13:\quad n^2-1=168=2^3\times3\times7,\ \text{trois distincts (2,3,7).}
\]
\[
\bullet\ n=14:\quad n^2-1=195=3\times5\times13,\ \text{trois distincts (3,5,13).}
\]

Les trois premiers $n$ qui conviennent : 
\[
n=11\ (\text{avec }120=2^3\times3\times5),\quad
n=13\ (\text{avec }168=2^3\times3\times7),\quad
n=14\ (\text{avec }195=3\times5\times13).
\]

\textbf{Réponse :} $\boxed{11,\ 13,\ \text{et }14.}$

\end{enumerate}
\end{solutionbox}

%%%%%%%%%%%%%%%%%%%%%%%%%%%%%%%%%%%%%%%%%%%%%%%%%%%%%%%%%%%%%%%%%%%%%%
% Exercice 16
%%%%%%%%%%%%%%%%%%%%%%%%%%%%%%%%%%%%%%%%%%%%%%%%%%%%%%%%%%%%%%%%%%%%%%
\begin{exercice}[title=Exercice 16 $\star\star$ -- {[Calculer, Raisonner]}]
Montrer que pour tout entier premier $p$ et pour tous entiers $x$ et $y$, on a :
$$
(x + y)^p \equiv x^p + y^p \pmod{p}.
$$
\end{exercice}

\begin{solutionbox}{Solution de l'exercice 16}

\textbf{1. Rappel du binôme de Newton mod $p$}

On sait que 
\[
(x+y)^p = \sum_{k=0}^{p} \binom{p}{k} x^{p-k} y^k.
\]
On montre en même temps (ou on admet) que si $p$ est premier, $\binom{p}{k}\equiv 0\pmod{p}$ pour $1\le k\le p-1.$

\textbf{2. Raisonnement sur $\binom{p}{k}$}

Pour $p$ premier et $1\le k\le p-1,$  
\[
\binom{p}{k} = \frac{p!}{k!(p-k)!}.
\]
Le numérateur contient un facteur $p,$ alors que $k!(p-k)!$ n’en contient pas (car $k<p,$ etc.). Du coup, $p$ divise $\binom{p}{k},$ donc 
\[
\binom{p}{k}\equiv0 \pmod{p}.
\]

\textbf{3. Conclusion sur $(x+y)^p$ mod $p$}

Donc dans la somme de Newton :
\[
(x+y)^p 
\equiv \binom{p}{0}x^p + \binom{p}{p}y^p
\pmod{p}.
\]
Or $\binom{p}{0}=1,$ $\binom{p}{p}=1.$ D’où 
\[
(x+y)^p \equiv x^p + y^p\pmod{p}.
\]

\end{solutionbox}

%%%%%%%%%%%%%%%%%%%%%%%%%%%%%%%%%%%%%%%%%%%%%%%%%%%%%%%%%%%%%%%%%%%%%%
% Exercice 17
%%%%%%%%%%%%%%%%%%%%%%%%%%%%%%%%%%%%%%%%%%%%%%%%%%%%%%%%%%%%%%%%%%%%%%
\begin{exercice}[title=Exercice 17 $\star\star$ -- [Chercher]]
Montrer que pour tout $n \le 10^9$, la décomposition de $n$ en produit de facteurs premiers fait apparaître moins de dix facteurs premiers distincts.
\end{exercice}

\begin{solutionbox}{Solution de l'exercice 17}

\textbf{Idée :} On veut prouver qu’un entier jusqu’à $10^9$ ne peut pas avoir 10 \emph{distincts} facteurs premiers.

\textbf{Raisonnement :}
\begin{enumerate}
  \item Si un nombre $n$ possédait 10 facteurs premiers distincts, notons-les $p_1<p_2<\dots <p_{10}.$  
  \item Le plus petit produit de 10 premiers distincts est 
    $$
    2\times 3\times 5\times7\times11\times13\times17\times19\times23\times29.
    $$
    Ce produit dépasse déjà largement $10^9$ (on peut estimer, $2\cdot3\cdot5\cdot7=210, \times11=2310,\times13=30030,\times17\approx 510510,\times19\approx9.7\times10^6,\times23>2.2\times10^8,\times29>6.4\times10^9...).$ Clairement $>10^9.$
  \item Ainsi, impossible d’avoir 10 distincts facteurs premiers sous la borne $10^9.$
\end{enumerate}

\textbf{Conclusion :} Tout $n\le 10^9$ admet dans sa décomposition en facteurs premiers strictement moins de 10 facteurs premiers distincts.

\end{solutionbox}

\begin{exercice}[title=Exercice 18 $\star$ { [Calculer] }]{}{}
    Dans chaque cas, déterminer le reste de la division euclidienne de $n$ par $a$ :
    \begin{enumerate}
      \item $n = 21^0$ et $a = 11$.
      \item $n = 31^7$ et $a = 17$.
      \item $n = 5^{20}$ et $a = 29$.
      \item $n = 13^4$ et $a = 7$.
      \item $n = 15^{100}$ et $a = 97$.
    \end{enumerate}
    \end{exercice}
    
    \begin{solutionbox}
    \textbf{Corrigé (sans utiliser la formule d’Euler, uniquement des réductions successives et le petit théorème de Fermat) :}
    
    \begin{itemize}
      \item \textbf{a)} $21^0 = 1$.  
      Le reste de la division de $1$ par $11$ est $1$.
    
      \item \textbf{b)} Calcul de $31^7 \bmod 17$.  
      Réduction d’abord : $31 \equiv 14 \pmod{17}$.  
      On doit donc évaluer $14^7 \bmod 17$.  
      $$14^2 = 196 \equiv 196 - 187 = 9 \pmod{17}.$$  
      $$14^3 \equiv 9 \times 14 = 126 \equiv 126 - 119 = 7 \pmod{17}.$$  
      $$14^4 \equiv 7 \times 14 = 98 \equiv 98 - 85 = 13 \pmod{17}.$$  
      $$14^5 \equiv 13 \times 14 = 182 \equiv 182 - 170 = 12 \pmod{17}.$$  
      $$14^6 \equiv 12 \times 14 = 168 \equiv 168 - 153 = 15 \pmod{17}.$$  
      $$14^7 \equiv 15 \times 14 = 210 \equiv 210 - 187 = 23 \equiv 6 \pmod{17}.$$  
      Le reste est donc $6$.
    
      \item \textbf{c)} $5^{20} \bmod 29$.  
      On réduit pas à pas (ou on utilise des puissances successives) :  
      \begin{align*}
        5^2 &= 25 \pmod{29},\\
        5^4 &= 25^2 = 625 \equiv 625 - 609 = 16 \pmod{29},\\
        5^8 &= 16^2 = 256 \equiv 256 - 232 = 24 \pmod{29},\\
        5^{16} &= 24^2 = 576 \equiv 576 - 551 = 25 \pmod{29}.
      \end{align*}
      Puis $20 = 16 + 4$. Donc :
      $$5^{20} = 5^{16} \times 5^4 \equiv 25 \times 16 = 400 \pmod{29}.$$  
      Or $400 - 377 = 23$ (car $29 \times 13 = 377$), donc $5^{20} \equiv 23 \pmod{29}$.  
      Reste $23$.
    
      \item \textbf{d)} $13^4 \bmod 7$.  
      $13 \equiv -1 \pmod{7}$, donc $13^2 \equiv (-1)^2 = 1 \pmod{7}$ et $13^4 \equiv 1^2 = 1 \pmod{7}$.  
      Reste $1$.
    
      \item \textbf{e)} $15^{100} \bmod 97$.  
      Ici, on peut utiliser le \emph{petit théorème de Fermat} (car $97$ est un nombre premier). Il dit : pour tout entier $a$ non divisible par $97$, $a^{96} \equiv 1 \pmod{97}$.  
      Comme $15$ n’est pas multiple de $97$, on a :  
      $$15^{96} \equiv 1 \pmod{97}.$$  
      On factorise :  
      $$15^{100} = 15^{96} \times 15^4 \equiv 1 \times 15^4 \pmod{97}.$$  
      Restent donc à calculer $15^4 \bmod 97$ :  
      $$15^2 = 225 \equiv 225 - 194 = 31 \pmod{97} \quad (\text{car } 97 \times 2 = 194).$$  
      $$15^4 \equiv 31^2 = 961 \equiv 961 - 873 = 88 \pmod{97} \quad (\text{car } 97 \times 9 = 873).$$  
      Donc $15^{100} \equiv 88 \pmod{97}$.  
      Reste $88$.
    \end{itemize}
    \end{solutionbox}
    
    %---------------------------------------------------------------------
    
    \begin{exercice}[title=Exercice 19 $\star\star$ { [Calculer] }]{}{}
    \begin{enumerate}
      \item Calculer $3^{24} \bmod 35$.
      \item Calculer $4^{20} \bmod 55$.
    \end{enumerate}
    \end{exercice}
    
    \begin{solutionbox}
    \textbf{Corrigé (sans Euler, via décomposition en facteurs premiers et le petit théorème de Fermat)} :
    
    \begin{enumerate}
      \item \textbf{$3^{24} \bmod 35$:}  
      On note que $35 = 5 \times 7$.  
      \begin{itemize}
        \item \textbf{Modulo 5 :} 
          Le petit théorème de Fermat dit que $3^4 \equiv 1 \pmod{5}$ (en effet, $3^2=9\equiv -1$, donc $3^4\equiv 1$).  
          Comme $24$ est un multiple de $4$, on a $3^{24} \equiv (3^4)^6 \equiv 1^6 = 1 \pmod{5}$.
        \item \textbf{Modulo 7 :} 
          De même, $3^6 \equiv 1 \pmod{7}$ (on peut le vérifier : $3^2=9\equiv 2,\ 3^3=27\equiv 6,\ 3^6=(3^3)^2\equiv 6^2=36\equiv 1\pmod7$).  
          Puisque $24$ est un multiple de $6$, $3^{24} = (3^6)^4 \equiv 1^4 = 1 \pmod{7}$.
      \end{itemize}
      Par le théorème chinois, $3^{24}\equiv 1$ modulo 5 et 7, donc $3^{24}\equiv 1 \pmod{35}$.  
      Reste $1$.
    
      \item \textbf{$4^{20} \bmod 55$:}  
      $55 = 5 \times 11$.  
      \begin{itemize}
        \item \textbf{Modulo 5 :} 
          $4 \equiv -1 \pmod{5}$, donc $4^{20} = (-1)^{20} = 1 \pmod{5}$.
        \item \textbf{Modulo 11 :} 
          On utilise encore le petit théorème de Fermat : $4^{10} \equiv 1 \pmod{11}$.  
          Or $20$ est $2\times10$, donc $4^{20}\equiv (4^{10})^2\equiv 1^2=1 \pmod{11}$.  
          (On peut aussi faire directement les calculs de petites puissances de 4 modulo 11.)
      \end{itemize}
      Ainsi $4^{20}\equiv 1$ modulo 5 et 11, donc $4^{20}\equiv 1 \pmod{55}$.  
      Reste $1$.
    \end{enumerate}
    \end{solutionbox}
    

%=====================================================================
% Exemple de mise en forme conforme aux instructions précédentes
% (exercice + solutionbox, math entre $...$ et $$...$$)
%=====================================================================

% \begin{exercice}[title=Exercice 18 $\star$ { [Calculer] }]{}{}
%     Dans chaque cas, déterminer le reste de la division euclidienne de $n$ par $a$ :
%     \begin{enumerate}
%       \item $n = 21^0$ et $a = 11$.
%       \item $n = 31^7$ et $a = 17$.
%       \item $n = 5^{20}$ et $a = 29$.
%       \item $n = 13^4$ et $a = 7$.
%       \item $n = 15^{100}$ et $a = 97$.
%     \end{enumerate}
%     \end{exercice}
%     
%     \begin{solutionbox}
%     \textbf{Corrigé :}
%     
%     \begin{itemize}
%       \item \textbf{a)} $21^0 = 1$. Le reste de 1 par 11 est 1.
%       
%       \item \textbf{b)} Calculer $31^7 \bmod 17$.  
%       D’abord, $31 \equiv 14 \pmod{17}$.  
%       On doit donc trouver $14^7 \bmod 17$.  
%       $$
%         14^2 = 196 \equiv 9 \pmod{17},\quad
%         14^3 \equiv 9 \times 14 = 126 \equiv 7 \pmod{17},
%       $$
%       $$
%         14^4 \equiv 7 \times 14 = 98 \equiv 13,\quad
%         14^5 \equiv 13 \times 14 = 182 \equiv 12,
%       $$
%       $$
%         14^6 \equiv 12 \times 14 = 168 \equiv 15,\quad
%         14^7 \equiv 15 \times 14 = 210 \equiv 6 \pmod{17}.
%       $$
%       D’où le reste 6.
%     
%       \item \textbf{c)} $5^{20} \bmod 29$.  
%       Méthode «puissance carrée» :
%       $$
%         5^2 = 25 \pmod{29},\quad
%         5^4 = 25^2 = 625 \equiv 16 \pmod{29},
%       $$
%       $$
%         5^8 = 16^2 = 256 \equiv 24,\quad
%         5^{16} = 24^2 = 576 \equiv 25 \pmod{29}.
%       $$
%       Puisque $20 = 16 + 4$,
%       $$
%         5^{20} = 5^{16}\times 5^4 \equiv 25 \times 16 = 400 \equiv 23 \pmod{29}.
%       $$
%       Reste 23.
%     
%       \item \textbf{d)} $13^4 \bmod 7$.  
%       $$
%         13 \equiv -1 \pmod{7} \quad\Longrightarrow\quad 13^2 \equiv 1,\quad 13^4 \equiv (13^2)^2 \equiv 1.
%       $$
%       Reste 1.
%     
%       \item \textbf{e)} $15^{100} \bmod 97$.  
%       $97$ est premier, donc $\varphi(97) = 96$ et $15^{96} \equiv 1 \pmod{97}$.  
%       On a 
%       $$
%         15^{100} = 15^{96}\times 15^4 \equiv 1 \times 15^4 \pmod{97}.
%       $$
%       Or
%       $$
%         15^2 = 225 \equiv 31,\quad
%         15^4 \equiv 31^2 = 961 \equiv 88 \pmod{97}.
%       $$
%       Donc le reste est 88.
%     \end{itemize}
%     \end{solutionbox}
%     
%     %---------------------------------------------------------------------
%     
%     \begin{exercice}[title=Exercice 19 $\star\star$ { [Calculer] }]{}{}
%     \begin{enumerate}
%       \item Calculer $3^{24} \bmod 35$.
%       \item Calculer $4^{20} \bmod 55$.
%     \end{enumerate}
%     \end{exercice}
%     
%     \begin{solutionbox}
%     \textbf{Corrigé :}
%     
%     \begin{enumerate}
%       \item \textbf{$3^{24} \bmod 35$:}  
%       On remarque $\gcd(3,35)=1$ et $\varphi(35)=24$.  
%       Par le théorème d’Euler, 
%       $$
%         3^{24} \equiv 1 \pmod{35}.
%       $$
%       Le reste est 1.
%     
%       \item \textbf{$4^{20} \bmod 55$:}  
%       On a $55 = 5 \times 11$, $\gcd(4,55)=1$.  
%       Méthode : on calcule séparément modulo 5 et modulo 11.
%       \begin{itemize}
%         \item \textbf{Mod 5 :} $4 \equiv -1 \pmod{5}$, donc
%           $$
%             4^{20} \equiv (-1)^{20} = 1 \pmod{5}.
%           $$
%         \item \textbf{Mod 11 :} 
%           $$
%             4^1=4,\quad
%             4^2=16 \equiv 5,\quad
%             4^3=20 \equiv 9,\quad
%             4^4=36 \equiv 3,\quad
%             4^5=12 \equiv 1 \pmod{11}.
%           $$
%           Donc $(4^5)^4 = 4^{20} \equiv 1^4 = 1 \pmod{11}$.
%       \end{itemize}
%       Par le théorème chinois, $4^{20} \equiv 1 \pmod{55}$.  
%       Le reste est 1.
%     \end{enumerate}
%     \end{solutionbox}
%     
%     %---------------------------------------------------------------------
    
    \begin{exercice}[title=Exercice 20 $\star\star$ { [Calculer] }]{}{}
    Déterminer le chiffre des unités de $3^{80}$ et de $7^{28}$.
    \end{exercice}
    
    \begin{solutionbox}
    \textbf{Corrigé :}
    
    \paragraph{Chiffre des unités de $3^{80}$.}
    La suite des unités de $3^n$ est périodique de période 4 :
    $$
    3^1 = 3,\quad
    3^2 = 9,\quad
    3^3 = 27 \ (\text{unité } 7),\quad
    3^4 = 81 \ (\text{unité } 1),
    $$
    et ça recommence (3, 9, 7, 1, ...).  
    Puisque $80$ est multiple de 4, le chiffre des unités de $3^{80}$ est 1.
    
    \paragraph{Chiffre des unités de $7^{28}$.}
    La suite des unités de $7^n$ (7, 9, 3, 1, ...) a également une période 4.  
    Comme $28$ est multiple de 4, le chiffre des unités de $7^{28}$ est 1.
    \end{solutionbox}
    


\end{document}
